% Copyright (c) 2017 by Daniel R. Driver.
% !TEX encoding = UTF-8 Unicode
% !TEX TS-program = XeLaTeX

\documentclass[titlepage]{article}

% This document presumes a file structure and set of inputs that are
% available at: git@github.com:danieldriver/syllabi.git

\newcommand\policy{../policy}
\newcommand\incl{../includes}
\ProvidesFile{variables.tex}[2022/09/11 v2.3 -- Syllabus variables]

\usepackage{xspace} % make manual spaces (like \mycmd\ ) unnecessary
\usepackage{xifthen} % provides \isempty test

% variables for internal use
\newcommand\prof{}
\newcommand\pdegree{}
\newcommand\pphone{}
\newcommand\pemail{}
\newcommand\poffice{}
\newcommand\phours{}
%
\newcommand\ccode{}
\newcommand\ctitle{}
\newcommand\cseries{}
\newcommand\cversion{}
\newcommand\csemester{}
\newcommand\cmeetson{}
\newcommand\cmeetsat{}
\newcommand\cmeetsin{}
\newcommand\cwebsite{}
\newcommand\cdescrip{}
\newcommand\cprereqs{}
\newcommand\edobject{}

% in case of fully online, BTh, or undergraduate courses - https://tex.stackexchange.com/a/5896
\newif\ifonline
\newcommand\Int[2]{\ifonline#1\else#2\fi}
\newif\ifbth
\newcommand\bth[2]{\ifbth#1\else#2\fi}
\newif\ifUG
\newcommand\UG[2]{\ifUG#1\else#2\fi}

% commands for setting variables in the preamble
\newcommand\professor[2][PhD]{
  \renewcommand\pdegree{#1\xspace}
  \renewcommand\prof{#2\xspace}}
\newcommand\phone[1]{
  \renewcommand\pphone{\addfontfeatures{Numbers=Monospaced}#1\xspace}}
\newcommand\email[1]{
  \renewcommand\pemail{\href{mailto:#1}{#1}\xspace}}
\newcommand\officehours[2][AST Library/Teams]{
  \renewcommand\poffice{#1\xspace}
  \renewcommand\phours{\href{https://outlook.office.com/bookwithme/user/b97b147a2e5c43b4a0e200c2cb630b77@astheology.ns.ca?anonymous}{#2}\xspace}%
  }
%
\newcommand\coursecode[2][1.0]{
  \renewcommand\cversion{#1\xspace}
  \renewcommand\ccode{#2\xspace}}
\newcommand\coursetitle[2][]{
  \ifthenelse{\isempty{#1}}%
    {}% do nothing if #1 is empty, else:
    {\renewcommand\cseries{#1\\[1ex]}}
  \renewcommand\ctitle{#2\xspace}}
\newcommand\semester[1]{
  \renewcommand\csemester{#1\xspace}}
\newcommand\meets[3]{
  \newcommand\AM{\textsc{am}}
  \newcommand\PM{\textsc{pm}}
  \renewcommand\cmeetson{#1\xspace}
  \renewcommand\cmeetsat{\Int{From 9:00 \AM}{#2}\xspace}
  \renewcommand\cmeetsin{\Int{\href{https://smu.brightspace.com/d2l/login}{Brightspace}}{#3}\xspace}}
\newcommand\website[2][]{
  \renewcommand\cwebsite{
  \ifthenelse{\isempty{#1}}%
    {\href{https://#2}}% copy url if #1 is empty, else
    {\href{#1}}% use specified url and print #2
  {#2}\xspace}}
\newcommand\cdescription[2][BF 1001; RM 1000 or GTRS 6000]{
  \renewcommand\cprereqs{#1}
  \renewcommand\cdescrip{#2\par}}
\newcommand\objectives[1]{
  \renewcommand\edobject{#1\par}}


\coursecode[1.1.0]{HB/ST 3001}
\coursetitle[Hebrew Bible/Systematic Theology]{Christian Doctrine and\\[.8ex] the Old Testament}

% Taught as RLGS 4483 (a remote antecedent) at Tyndale UC in:
%   - Winter 2009 ("Biblical Theological Themes: The Old Testament as Christian Scripture")
% Taught as HB/ST 3001 (largely a new course) at AST in:
%   - Summer 2017 (Session 3: July 17-21 and 24-28)

\professor{Daniel R. Driver}
\phone{902-425-7051}
\email{ddriver@astheology.ns.ca}
\officehours{Daily, 12:15--12:45 \PM}

\semester{Summer Term 2017}
\meets{Mon–Fri, July 17–28}% \meets{on}{at}{in}
      {9:00--12:00 \AM}
      {Classroom 2}
\website{danieldriver.com/courses/}
\cdescription{% copy from the current Academic Calendar
	What is the relationship of scripture and theology? Some answers to
	this question start with the Bible, while others begin with
	Christian doctrine, while still others call for a dialectical
	movement between canon and creed. However it is approached, the
	question presents a number of challenges.

	To illustrate the problem in general, and to outline some possible
	solutions, this course highlights a variety of current approaches to
	the Old Testament as the scripture of the Church. Students will
	encounter work by an array of biblical and theological specialists
	who transgress the conventional limits of their disciplines. They
	will study the views of Catholics, Protestants, Jews, and those who
	disavow religious affiliation. Because the Hebrew Bible is Jewish
	scripture in the first instance, students will consider why some
	Jews are not interested in biblical theology, and what it means for
	the Church to adhere to the Scriptures of Israel. Specific doctrines
	under consideration include: creation, sin, election, incarnation
	(the body of God), divine impassibility, Christology, Mariology, law
	and gospel, death and resurrection, and the afterlife.
}% end of course description
\objectives{% review Bloom's taxonomy: http://www.celt.iastate.edu/teaching/effective-teaching-practices/revised-blooms-taxonomy
            % or try this new one at Rice: http://cte.rice.edu/blogarchive/2017/2/23/taxonomy
	By the end of this course students should
		develop knowledge of some ways different churches -- include the RCC, ACC and UCC --
		approach the Bible in general, and the Old Testament in particular, as scripture.
	They should become familiar with
		a few major narratives about the relationship of Bible and Theology, and with
		different theories about how to frame knowledge of God (in terms apophatic and kataphatic, for example).
	They should be able to analyze and evaluate doctrinal claims
		by biblical scholars and theologians who use the OT as a core theological resource.
	Students should also develop curiosity about the role of Hebrew scripture in
		communities beyond the church, including in the academy and especially the synagogue.
		Their creativity and confidence in expressing their own ideas about the Scriptures of Israel should grow,
		even as they learn to do so with appropriate humility and charity.
	Finally, students should experience a process of discovery in dialogue with
		biblical traditions ancient and strange as they come to terms with their own ways
		of thinking and speaking about the God of Israel.
}% end of learning objectives

\ProvidesFile{preamble.tex}[2014/09/02 v1.1 -- Syllabus preamble]

% basic typography
\usepackage{fontspec,realscripts}
\defaultfontfeatures{Ligatures=TeX}
\setmainfont[%
  Numbers={Proportional,OldStyle},
  SmallCapsFeatures={RawFeature=+c2sc;+smcp,LetterSpace=1}]{Scala OT}
\newfontfamily\lining[Numbers={Lining}]{Scala OT}
\setsansfont[%
  SmallCapsFeatures={RawFeature=+c2sc;+smcp,LetterSpace=1}]{Scala Sans Pro}
\setmonofont[Scale=MatchLowercase]{Menlo}
\usepackage{sectsty}
\allsectionsfont{\sffamily\scshape}
\frenchspacing
\setlength{\emergencystretch}{3em} % prevent overfull lines

% custom font size and leading
\renewcommand\tiny{\fontsize{6}{9}\selectfont}
\renewcommand\scriptsize{\fontsize{7}{10}\selectfont}
\renewcommand\footnotesize{\fontsize{8}{10}\selectfont}
\renewcommand\small{\fontsize{9}{11}\selectfont}
\renewcommand\normalsize{\fontsize{10}{12}\selectfont}% base size
\renewcommand\large{\fontsize{11}{13}\selectfont}
\renewcommand\Large{\fontsize{12}{14}\selectfont}
\renewcommand\LARGE{\fontsize{16}{18}\selectfont}% "course syllabus \\ semester" benefits from more lead
\renewcommand\huge{\fontsize{20}{22}\selectfont}
\renewcommand\Huge{\fontsize{24}{26}\selectfont}

% layout packages: page, logo, tables
\usepackage[scale={0.6,0.8},
            xetex]{geometry}
\usepackage{graphicx}
\usepackage{array}     % allow insertions of column styling with >{}
\usepackage{booktabs}  % elegant horizontal rules in tables

% custom macros for a session count in the schedule of readings
\newcounter{session}
\newcounter{columns}
\newcounter{spancolumns}
\newcounter{midcolumns}
\newcounter{courseunit}
\newcommand\setcolumncount[2][0]{ % optionally set count to other than 0,
  \setcounter{session}{#1}        % e.g. to -1, or to a standing count
  \setcounter{columns}{#2}
  \setcounter{spancolumns}{#2-1}
  \setcounter{midcolumns}{#2-2}}
\newcommand\sessioncount{\stepcounter{session}\arabic{session}}
\newcommand\sessionskip[1]{\multicolumn{1}{@{}r@{ }}{#1}}
\newcommand\unit[1]{\multicolumn{\thecolumns}{c}{%
  \scshape\stepcounter{courseunit}\roman{courseunit}. #1}}
\newcommand\noclass[1]{\sessionskip{---.}&\multicolumn{\themidcolumns}{@{}l}{\itshape No Class: #1}}
\newcommand\rarr{\Large\char"EA31\hspace*{0.5pt}} % pointer finger w/ reminder tie in Scala OT
\newcommand\reminder[2]{\sessionskip{\rarr}&\multicolumn{\thespancolumns}{@{}l}{#1 \hfill #2}}

% color to match Tyndale's branding
\usepackage[usenames]{xcolor}
% predefined: black, white, red, green, blue, cyan, magenta, yellow
\definecolor{TyndaleURLs}{HTML}{0062A0} % links on tyndale.ca
\definecolor{TyndaleBlue}{cmyk}{1,1,0,.32}
\definecolor{TyndaleGold}{cmyk}{0,.27,1,0}
\definecolor{TyndaleRed}{cmyk}{0,1,.99,.04}
\definecolor{TyndaleBlack}{cmyk}{0,0,0,1}
\definecolor{TyndaleGreen}{cmyk}{.45,0,1,.24}
\definecolor{TyndaleOrange}{cmyk}{0,.79,1,0}
\definecolor{TyndaleAqua}{cmyk}{.47,0,.24,0}
\definecolor{TyndaleYellow}{cmyk}{.03,.03,.35,0}

% metadata (assumes a host of definitions are made in the main file)
\usepackage[setpagesize=false,     % leave this to geometry
            hyperfootnotes=false,  % fragile and distracting
            xetex]{hyperref}
\hypersetup{breaklinks=true,       % allow link text to break across lines
            colorlinks=true,       % colorlinks resets pdfborder to 0 0 0
            urlcolor=TyndaleURLs,  % for external links
            linkcolor=TyndaleRed,  % for normal internal links
            citecolor=TyndaleGold, % for bibliographical citations in text
            pdfauthor={\prof},
            pdftitle={\ccode: \ctitle},
            pdfsubject={Tyndale UC, \csemester},
            pdfcreator={github.com/danieldriver/syllabus}}
\urlstyle{same}                    % don't use monospace font for urls

% custom footlines
\usepackage{fancyhdr}
\pagestyle{fancy} % turn it on
\fancyhf{}        % reset everything
\renewcommand{\headrulewidth}{0pt} % remove header line as well
\lfoot{\sffamily\scshape\footnotesize\ctitle, \csemester}
\rfoot{\sffamily\scshape\footnotesize\prof\quad\thepage}

% language support: polyglossia needs to go last
\usepackage{polyglossia}
\setdefaultlanguage{english}
\setotherlanguage{hebrew}
\newfontfamily\hebrewfont[Script=Hebrew]{SBL BibLit}

% protect positioning of margin table in policy/grades -- needs to come after polyglossia
\usepackage{marginfix}

% gratuitous with custom title page, but useful as a fallback
\title{\ccode: \ctitle}
\author{\professor}
\date{\semester}


\begin{document}
\ProvidesFile{title.tex}[2013/09/06 v1.0 -- Syllabus title page]

\begin{titlepage}
  \begin{center}

    \begin{minipage}{\textwidth}
      \parbox[t]{0.5\textwidth}{
        \mbox{}\\[-\baselineskip] % dummy line to align parboxes
        \includegraphics[width=0.5\textwidth]{.syllabus/includes/TyndaleUC}}
      \hfill
      \parbox[t]{0.4\textwidth}{
        \raggedleft\LARGE\sffamily Course Syllabus\\
        \semester}
    \end{minipage}

    \vfill

    {\sffamily\textsc{\LARGE\MakeLowercase\coursecode}\\[2ex]
      \bfseries\itshape\Huge\coursetitle}

    \vfill

    \begin{tabular}{>{\bfseries}rl>{\bfseries}rl}
      \toprule
      Instructor & \professor                   & Course  & \liningnums{\coursecode}~\texttt{(v\version)} \\
      \midrule
      Phone      & \texttt{\phone}              & Meets   & \meeton                                       \\
      Email      & \href{mailto:\email}{\email} & Time    & \meetat                                       \\
      Office     & \officenumber                & Room    & \meetin                                       \\
      Hours      & \officehours                 & Website & \href{http://\website}{\website}              \\
      \bottomrule
    \end{tabular}

    \vfill

    \begin{description}
      \item[Commuter Hotline]
        Class cancellations due to inclement weather or illness will
        be announced on the commuter hotline at \texttt{416.226.6620
        x2187}. Alternately, weather cancellation information is posted
        at \href{http://tyndale.ca/weather}{tyndale.ca/weather}.
      \item[MyTyndale.ca]
        This course may have materials stored on its website, such as
        handouts or readings that may be needed in order to complete
        assignments. Students are responsible for checking these course
        pages on a regular basis. Here, too, students are able to view
        their grades throughout the semester. For more information see
        Section~\ref{mytyndale}, below.
      \item[Mail]
        Students are responsible for information communicated through
        their campus mailboxes and student e-mail accounts. A mailbox
        directory hangs beside the mailboxes. For more information
        contact the Registrar's office.
    \end{description}

  \end{center}

  \section{Course Description}
  \label{description}

  \emph{From the Academic Calendar:} \cdescription

\end{titlepage}
\setcounter{page}{2} % count the title page as page 1


\section{Learning Objectives}
\label{objectives}
\edobject

\section{Required Texts \& Materials}
\label{texts}

The following texts are required. Students are encouraged to purchase
their own copies. Library copies that are not already reference works
will be placed on a 1-hour reserve.

\begingroup
\renewcommand{\section}[2]{}% temporarily remove the section heading
\begin{thebibliography}{MacDonald}% use the longest item in the bibliography

	\bibitem[NRSV]{nrsv} M.\,D. Coogan, ed.
    \emph{New Oxford Annotated Bible with Apocrypha: NRSV}. 4th ed.
    Oxford / New York: Oxford University Press, 2010.
%    \textsc{isbn} 978-0195289602.

	\bibitem[Anderson]{anderson} Gary A. Anderson.
	\emph{Christian Doctrine and the Old Testament: Theology in the Service of Biblical Exegesis}.
	Grand Rapids: Baker Academic, 2017.
%	\textsc{isbn} 978-0801098253.

	\bibitem[MacDonald]{macdonald} Neil B. MacDonald.
	\emph{Metaphysics and the God of Israel: Systematic Theology of the Old and New Testaments}.
	Grand Rapids: Baker Academic, 2006.
%	\textsc{isbn} 978-0801032431.
	For those who are unable to find copies of this book for resale, the library's reserve copy should suffice to meet the needs of the class.

%	\bibitem[tk]{tk} tk.
%	\emph{tk}.
%	tk: tk, 20tk.
%	\textsc{isbn} 978-tk.

\end{thebibliography}
\endgroup

An acceptable alternative study Bible is the NJPS: Adele Berlin and Marc
Zvi Brettler, eds., \emph{The Jewish Study Bible: Second Edition}
(Oxford / New York: Oxford University Press, 2014). A reference copy is
available in the library, and it is well worth consulting.

\section{Supplementary Texts}
\label{supplementary}

Supplementary readings will be recommended throughout the course. In
particular, students are asked to explore relevant church statements
about Christian scripture, particularly the Old Testament. Resources
from the founding churches of AST include the following.

\begin{enumerate}

\item RCC: Roman Catholic Church

	\begin{enumerate}
	\item \href{}{\emph{Divino Afflante Spiritu}, the encyclical from Pope Pius XII [\href{http://www.bc.edu/schools/stm/crossroads/resources/deathofjesus/intro/the_catholic_approachtothebible.html}{overview}] (1943)}
	\item \href{}{\emph{Dei Verbum}, Vatican II's Dogmatic Constitution on Divine Revelation (1965)}
	\item PBC = Pontifical Biblical Commission: \emph{L'interprétation de la Bible dans l'Église}, \href{http://www.catholic-resources.org/ChurchDocs/PBC_Interp.htm}{The Interpretation of the Bible in the Church} (April 15, 1993)
	\item PBC: \emph{Le peuple juif et ses Saintes Écritures dans la Bible chrétienne}, \href{http://www.vatican.va/roman_curia/congregations/cfaith/pcb_documents/rc_con_cfaith_doc_20020212_popolo-ebraico_en.html}{The Jewish People and Their Sacred Scriptures in the Christian Bible} (May 24, 2001)
	\end{enumerate}

\item ACC: Anglican Church of Canada (one of 38 Provinces in the Anglican Communion)

	\begin{enumerate}
	\item Articles VI–VII of \href{https://www.churchofengland.org/prayer-worship/worship/book-of-common-prayer/articles-of-religion.aspx}{The Thirty-nine Articles of Religion} (1562)
	\item \href{http://anglicansonline.org/basics/Chicago_Lambeth.html}{The Chicago-Lambeth Quadrilateral} (1886, 1888)
	\item Paragraphs 52–62 of \href{http://www.anglicancommunion.org/resources/document-library.aspx?author=The+Windsor+Process&language=English}{The Windsor Report} [\href{http://www.anglicancommunion.org/media/68225/windsor2004full.pdf}{PDF}] (18 October 2004)
	\item \href{http://www.anglicancommunion.org/media/98131/Final-Report-for-the-web.pdf}{Bible in the Life of the Church (BILC) Project Report} (2012)
	\end{enumerate}

\item UCC: United Church of Canada

	\begin{enumerate}
	\item The UCC requires its ministers to be in ``\href{http://www.united-church.ca/community-faith/welcome-united-church-canada/what-we-believe}{essential agreement}'' with four doctrinal statements that are ``\href{http://www.united-church.ca/community-faith/welcome-united-church-canada/faith-statements}{subordinate to the primacy of scripture}'':
		\begin{enumerate}
			\item\href{http://www.united-church.ca/community-faith/welcome-united-church-canada/twenty-articles-doctrine-1925}{Twenty Articles of Doctrine} (1925)
			\item\href{http://www.united-church.ca/community-faith/welcome-united-church-canada/statement-faith-1940}{A Statement of Faith} (1940)
			\item\href{http://www.united-church.ca/community-faith/welcome-united-church-canada/new-creed}{A New Creed} (1968)
			\item\href{http://www.united-church.ca/community-faith/welcome-united-church-canada/song-faith}{A Song of Faith} (2006)
		\end{enumerate}
	\item \href{https://ecumenism.net/archive/docu/1992_ucc_authority_interpretation_scripture.pdf}{\emph{The Authority and Interpretation of Scripture: A Statement of the United Church of Canada}} (Toronto: UC Publishing House, 1992) [BS 480 U65 1992 -- \textbf{on reserve}]
	\item See also R.\,C. Fennell's series “How Does the United Church Interpret the Bible?” in \href{http://touchstonecanada.ca}{\emph{Touchstone}}, a Canadian journal of theology and heritage:
		\begin{enumerate}
			\item\href{http://touchstonecanada.ca/wp-content/uploads/2013/08/Article2_2008-051.pdf}{Part I, 1904–1940s: Tradition and Resistance} (26.2 [2008]: 13--24)
			\item\href{http://touchstonecanada.ca/wp-content/uploads/2013/08/3rdArticle_2008-09.pdf}{Part II, 1950s–1990s: Tradition and Liberation} (26.3 [2008]: 31--42)
			\item\href{http://touchstonecanada.ca/wp-content/uploads/2013/08/Touchstone_May_2011.pdf}{Part III: A Song of Faith} (29.2 [2011]: 21--29)
		\end{enumerate}
	\end{enumerate}

\end{enumerate}


\section{Course Outline}
\label{outline}

We will adhere to the schedule in \autoref{schedule} as closely as
possible, though the professor reserves the right to adjust it to suit
the needs of the class.

\setcolumncount{6}% set up \sessioncount, \unit{}, \noclass{}, and \reminder{memo}{date} macros
\begin{table}[htbp]% set to `p' to put the schedule on its own page
  \centering
  \begin{tabular}{>{\sessioncount.}r@{ }lcccr}% make sure the column config agrees with \setcolumncount
	\toprule
	\sessionskip{\textbf{\S}.}&\textbf{Topic}&\textbf{Anderson}&\textbf{MacDonald}&\textbf{OT Text}&\textbf{Date}\\
	\midrule

		& Approaches          & 1    & 1    & 1 Sam 4–6      & 17 July \\
		& Impassibility       & 2    & 2--4 & Jonah 1--4     & 18 July \\
		& Creation            & 3    & 5--6 & Gen 1--3       & 19 July \\
	\reminder{Articles of Faith \textbf{due} at the start of class four}{} \\
		& Original Sin        & 4    &      & Exod 32--34    & 20 July \\
		& Election            & 5    & 7--8 & Gen 12, 15, 17 & 21 July \\
	\noclass{Weekend}                                        & 22--23 July \\
		& Christology         & 6    & 9    & 1 Kings 8      & 24 July \\
		& Mariology           & 7    &      & 1 Sam 1--2     & 25 July \\
		& Faith \& Works      & 8--9 & 10   & Tobit 1--4     & 26 July \\
		& Atonement           & 10   & 11   & Tobit 12--14   & 27 July \\
	\reminder{Paper is \textbf{due} at the start of class ten}{}  \\
		& Accordance          &      &      & 2 Kings 17     & 28 July \\
	\reminder{End of Term: Final marks are due to the Registrar}{7 Aug.} \\

	\bottomrule
  \end{tabular}
  \caption{Schedule of Readings}
  \label{schedule}
\end{table}

See the AST website for a list of other \href{http://www.astheology.ns.ca/students/academic-dates.html}{important dates}.

\section{Evaluation and Grade Structure}
\label{evaluation}

\begin{enumerate}% 250*7+250*10

	\item Before each class, complete the assigned reading and write a
	one page (250 word) \textbf{reading response}. Use the exercise to
	elaborate, reflect on, test, challenge, or extend an idea from the
	reading that you find noteworthy. Type your responses, write in
	complete sentences, and combine those sentences into one or two
	healthy paragraphs. Due at the start of \textbf{each class}, unless
	it is the first class or another graded assignment is due. Students
	who submit fewer than 7 responses cannot earn an “A” for the course.

	\item On a single sheet of paper, annotate a list of between 7 and
	12 essential \textbf{articles of faith} that Christians historically
	derive from Old Testament scripture. For guidance on what counts as
	“essential” you may consult various catechisms (Catholic, Lutheran),
	creeds (Apostles', Nicene, New), and other official documents
	(Articles, Encyclicals, Statements). In your annotations, cite OT
	texts that seem to speak to each of the doctrines you select,
	whether for or against or both. Due at the start of \textbf{class
	four}.

	\item A \textbf{final paper} allows you to explore one Christian
	Doctrine of your choice in terms of its relationship to Old
	Testament scripture. The paper must have a clear thesis and be 2,500
	words (ten pages) long. It should interact with the assigned reading
	and cite at least seven secondary sources in total. Due at the start
	of the \textbf{last class}.

\end{enumerate}

Note that this course is an intensive. Its format is highly condensed.
In a two week period students are expected to attend 30 hours of class
and to spend a further 50 hours \emph{or more} in independent study.
Students are therefore encouraged not to have other major commitments
during the course session so that they can have the best chance of
success.

The breakdown for the course's total work is shown in
\autoref{grade-dist}.

\begin{table}[htbp]
  \centering
  {\lining
  \begin{tabular}{lr}
    \toprule
    Reading Responses  & 40\% \\
    Annotated Articles & 20\% \\
    Final Paper        & 40\% \\
    \bottomrule
  \end{tabular}}
  \caption{Distribution of Grades}
  \label{grade-dist}
\end{table}

\ProvidesFile{grades.tex}[2023/09/06 v2.5 -- Course policy]

AST's \href{https://www.astheology.ns.ca/students/resources.html}{Academic
Calendar} provides guidelines and criteria for academic assessment. Marks are
assigned by letter grade using these benchmarks. Note that graduate and
undergraduate programs are currently assessed on a different scales at AST (see
\autoref{grade-syst}). For graduate students, a grade of C is deemed
``Unsatisfactory,'' whereas for undergraduate (BTh) students, a C is
``Satisfactory'' and the grade of D is possible for ``Marginal'' performance.
Other grades have different definitions, too.

\begin{table}[htbp]
 \centering
 {\lining\footnotesize
 \begin{tabular}{lll|ll}
   \toprule
               & \multicolumn{2}{c|}{Graduate Courses}   & \multicolumn{2}{c}{Undergraduate Courses} \\
               \cline{2-5}
   A+          & 94--100    & Exceptional               & 90--100    & Excellent                 \\
   A           & 87--93     & Outstanding               & 85--89     & Excellent                 \\
   A\char"2212 & 80--86     & Excellent                 & 80--84     & Excellent                 \\ [1ex]
   B+          & 77--79     & Good                      & 75--79     & Good                      \\
   B           & 73--76     & Acceptable                & 70--74     & Good                      \\
   B\char"2212 & 70--72     & Marginally Acceptable     &            &                           \\ [1ex]
   C           & 60--69     & Unsatisfactory            & 60--69     & Satisfactory              \\
   D           &            &                           & 50--59     & Marginal                  \\
   F           & 0--59      & Failure                   & 0--49      & Failure                   \\
   FP          & 0          & Failure due to Plagiarism & 0          & Failure due to Plagiarism \\
   \bottomrule
 \end{tabular}}
 \caption{Outline of Grades at AST}
 \label{grade-syst}
\end{table}

\ProvidesFile{other.tex}[2018/0/03 v2.4 -- Course policy]

\section{Other Course Policy}
\label{policy}

Late work will not be accepted, except in genuinely extenuating
circumstances. Students must submit something before the deadline if
they wish to receive credit. Unless I state otherwise, written
assignments are to be uploaded by 11:59 \PM\ on the date indicated.

Essay submissions must be typewritten, double-spaced, and formatted as
PDFs. They should be free from error. In this course they should follow
either \emph{The SBL Handbook of Style} (refer to the online
\href{https://www.sbl-site.org/assets/pdfs/pubs/SBLHSsupp2015-02.pdf}{
Student Supplement}), or \emph{The Chicago Manual of Style}, per AST's
policy for major papers. As a reminder, AST also upholds an Inclusive
Language Policy.

\href{http://www.eerdmans.com/Pages/Item/59043/Commentary-Statement.aspx}{Plagiarism},
if \href{https://www.theguardian.com/world/2013/feb/09/german-education-minister-quits-phd-plagiarism}{detected},
will result in failure of the course.

Students should request permission to record a class or lecture. If
permission is granted, or if recordings are provided (as in the case of
an online course), I stipulate that all recordings be for personal use
only. They may not be shared or distributed.

If you have abilities or disabilities that require modifications to the
assessment process or other aspects of this course, please advise the
course instructor as soon as possible.

Finally, I encourage the conscientious use of laptops, tablets, and
other technology in my classes. In classroom settings, realize that, as
\href{http://dx.doi.org/10.1016/j.compedu.2012.10.003}{cognitive
psychologists have demonstrated}, ``laptop multitasking hinders
classroom learning for both users and nearby peers.'' Do your part to
foster an environment of open dialogue by honouring the presence of your
classmates. In online settings, consider both the physical environment
in which you choose to work and the virtual environment that you help
create through your participation in various forums. Let your engagement
in this course be marked by rigour and charity alike.

\end{document}